\chapter{Teorema Pythagoras}\label{ch:g8:pythagoras}

Teorema yang paling terkenal dalam seluruh matematika menghubungkan
ketiga sisi \omterm{def:g6:shapes:triangles}{segitiga siku-siku}. Dengannya, panjang yang tidak dapat
diukur langsung penggaris mana pun menjadi terhitung: diagonal, jarak,
tinggi. Ini juga perjumpaan besar kita yang pertama dengan bedanya
sebuah teorema dan konversnya.

\section{Teoremanya}

\begin{theorem}[Pythagoras]\label{thm:g8:pythagoras:direct}
Kalau segitiga $ABC$ siku-siku di $A$, maka
\[
BC^2 = AB^2 + AC^2 :
\]
kuadrat sisi miringnya (yaitu sisi di seberang \omterm{def:g3:shapes:rightangle}{sudut siku-sikunya})
sama dengan \omterm{def:g1:addition:def}{jumlah} kuadrat kedua sisi lainnya.
\end{theorem}

\begin{proof}[Gagasan buktinya]
Empat salinan segitiganya, yang disusun di dalam persegi besar bersisi
$AB + AC$, menyisakan persegi miring yang terbangun pada sisi
miringnya tanpa tertutupi: membandingkan luasnya memberi kesamaannya.
Perhitungan rincinya diusulkan sebagai \cref{exo:g8:pythagoras:11};
teoremanya juga menyusul dari \omterm{def:g2:mult:def}{hasil kali} skalar, yang dipelajari di
buku Sekolah Menengah.
\end{proof}

\begin{omfigure}
\begin{tikzpicture}[scale=0.58]
  \coordinate (A) at (0,0);
  \coordinate (B) at (4,0);
  \coordinate (C) at (0,3);
  \draw[thick, fill=omDef!12] (A) -- (B) -- (C) -- cycle;
  \draw[thin] (0.35,0) -- (0.35,0.35) -- (0,0.35);
  \node[below left, font=\small] at (A) {$A$};
  \node[below right, font=\small] at (B) {$B$};
  \node[above left, font=\small] at (C) {$C$};
  % persegi pada sisinya
  \draw[omDef, thick, fill=omDef!8] (0,0) rectangle (4,-4);
  \node[omDef, font=\small] at (2,-2) {$AB^2 = 16$};
  \draw[omMeth, thick, fill=omMeth!8] (0,0) rectangle (-3,3);
  \node[omMeth, font=\small] at (-1.5,1.5) {$AC^2 = 9$};
  \draw[omThm, thick, fill=omThm!8] (4,0) -- (7,4) -- (3,7) -- (0,3)
    -- cycle;
  \node[omThm, font=\small] at (3.5,3.5) {$BC^2 = 25$};
\end{tikzpicture}

{\small Teoremanya sebagai pernyataan tentang luas: persegi yang
terbangun pada sisi miringnya ($25$) punya luas yang persis sama
dengan kedua persegi lainnya bersama-sama ($16 + 9$). Di sini sisinya
$3$, $4$, $5$.}
\end{omfigure}

\begin{method}[Menghitung sebuah panjang]\label{met:g8:pythagoras:compute}
Pada \omterm{def:g6:shapes:triangles}{segitiga siku-siku} yang dua sisinya diketahui:
\begin{enumerate}
  \item sebutkan sisi miringnya (yang di seberang \omterm{def:g3:shapes:rightangle}{sudut siku-sikunya}
        --- selalu sisi yang terpanjang);
  \item tulislah kesamaan \cref{thm:g8:pythagoras:direct} dengan nilai
        yang diketahui;
  \item carilah kuadrat yang hilang: \emph{jumlahkan} kuadratnya kalau
        sisi miringnya belum diketahui, \emph{kurangkan} kalau sisi
        siku-sikunya yang belum diketahui;
  \item ambillah akar kuadratnya, secara persis atau dengan
        kalkulator, lalu periksalah bahwa sisi miringnya keluar
        sebagai yang terpanjang.
\end{enumerate}
\end{method}

\begin{example}[Mencari sisi miringnya]\label{ex:g8:pythagoras:hypotenuse}
Sebuah \omterm{def:g6:shapes:triangles}{segitiga siku-siku} punya sisi siku-siku $6$ dan $8$. Sisi miringnya $c$:
\[
c^2 = 6^2 + 8^2 = 36 + 64 = 100,
\qquad
c = \sqrt{100} = 10 .
\]
(Bilangan $\sqrt{100}$, yaitu \emph{akar kuadrat} $100$, adalah
bilangan positif yang kuadratnya $100$; akar kuadrat memperoleh satu
bab penuh pada \cref{ch:g9:sqrt}.)
\end{example}

\begin{example}[Mencari sisi siku-siku]\label{ex:g8:pythagoras:leg}
Tangga sepanjang $5$ m bersandar pada tembok, kakinya $1.4$ m dari
tembok. Tinggi yang tercapai $h$: tangganya adalah sisi miringnya, jadi
\[
h^2 = 5^2 - 1.4^2 = 25 - 1.96 = 23.04,
\qquad
h = \sqrt{23.04} = 4.8 \text{ m}.
\]
Kurangkan, bukan jumlahkan: yang belum diketahui di sini adalah \emph{sisi siku-siku}.
\end{example}

\section{Konversnya}

\begin{theorem}[Konvers Pythagoras]\label{thm:g8:pythagoras:converse}
Kalau, pada segitiga $ABC$, sisinya memenuhi $BC^2 = AB^2 + AC^2$,
maka segitiganya siku-siku di $A$.
\end{theorem}
\admitted

\begin{method}[Menguji sudut siku-siku]\label{met:g8:pythagoras:test}
Diberikan ketiga sisi sebuah segitiga:
\begin{enumerate}
  \item hitunglah secara terpisah kuadrat sisi terpanjangnya, dan
        \omterm{def:g1:addition:def}{jumlah} kuadrat kedua sisi lainnya;
  \item kalau kedua hasilnya sama, segitiganya siku-siku (menurut
        konversnya), di \omterm{def:g5:solids:def}{titik sudut} yang di seberang sisi
        terpanjangnya;
  \item kalau keduanya berbeda, segitiganya tidak siku-siku ---
        karena teorema Pythagoras akan memaksa kesamaannya.
\end{enumerate}
\end{method}

\begin{example}\label{ex:g8:pythagoras:test}
Sisi $5$, $12$, $13$: kuadrat sisi terpanjangnya, $13^2 = 169$; \omterm{def:g1:addition:def}{jumlah}
kuadrat sisi lainnya, $5^2 + 12^2 = 25 + 144 = 169$. Sama: siku-siku
(di seberang sisi $13$).

Sisi $6$, $7$, $9$: $9^2 = 81$ tetapi $6^2 + 7^2 = 85 \neq 81$: bukan
\omterm{def:g6:shapes:triangles}{segitiga siku-siku} (ia ``hampir'' siku-siku --- bilangannya yang
memutuskan, bukan matanya).
\end{example}

\begin{remark}[Teorema melawan konversnya]
Teoremanya dan konversnya mengatakan hal yang berbeda: teoremanya
\emph{memakai} \omterm{def:g3:shapes:rightangle}{sudut siku-siku} untuk menghitung panjang; konversnya
memakai panjang untuk \emph{membuktikan} adanya \omterm{def:g3:shapes:rightangle}{sudut siku-siku}.
Tukang bangunan sudah memakai konversnya selama ribuan tahun: tali
bersimpul pada satuan $3$, $4$, $5$ yang ditarik tegang memberi \omterm{def:g3:shapes:rightangle}{sudut
siku-siku} yang sempurna.
\end{remark}

\begin{example}[Diagonal sebuah persegi]\label{ex:g8:pythagoras:diagonal}
Diagonal persegi bersisi $1$ adalah sisi miring \omterm{def:g6:shapes:triangles}{segitiga siku-siku}
dengan sisi siku-siku $1$ dan $1$:
\[
d^2 = 1^2 + 1^2 = 2 ,
\]
jadi $d = \sqrt2 \approx 1.414$ --- bilangan yang bukan \omterm{def:g6:fractions:def}{pecahan},
seperti yang akan dibuktikan buku Sekolah Menengah. Untuk persegi
bersisi $c$, diagonalnya adalah $c\sqrt2$.
\end{example}

\section{Latihan}

\begin{exercise}[$\star$]\label{exo:g8:pythagoras:1}
Pada tiap \omterm{def:g6:shapes:triangles}{segitiga siku-siku} ini, hitunglah sisi miringnya: sisi
siku-siku $9$ dan $12$; \quad $5$ dan $12$; \quad $8$ dan $15$.
\end{exercise}

\begin{exercise}[$\star$]\label{exo:g8:pythagoras:2}
Sebuah \omterm{def:g6:shapes:triangles}{segitiga siku-siku} punya sisi miring $25$ dan satu sisi
siku-siku $7$. Hitunglah sisi siku-siku yang lain.
\end{exercise}

\begin{exercise}[$\star$]\label{exo:g8:pythagoras:3}
Hitunglah diagonal persegi panjang bersisi $12$ cm dan $9$ cm; lalu
diagonal persegi bersisi $5$ cm (nilai persisnya dengan akar kuadrat,
lalu \omterm{def:g4:numbers:round}{dibulatkan} ke mm).
\end{exercise}

\begin{exercise}[$\star$]\label{exo:g8:pythagoras:4}
Segitiga yang mana yang siku-siku? Berilah alasan dengan
\cref{met:g8:pythagoras:test}:
\[
(20,\ 21,\ 29); \qquad
(7,\ 8,\ 11); \qquad
(1.5,\ 2,\ 2.5).
\]
\end{exercise}

\begin{exercise}[$\star$]\label{exo:g8:pythagoras:5}
Sebuah pintu gerbang diperkuat papan menyilang. Lebar gerbangnya
$3$ m dan tingginya $1.6$ m: berapa panjang papannya?
\end{exercise}

\begin{exercise}[$\star$]\label{exo:g8:pythagoras:6}
Sebuah \omterm{def:g6:shapes:triangles}{segitiga sama kaki} punya dua sisi $10$ cm dan alas $12$ cm.
Tingginya membelahnya menjadi dua \omterm{def:g6:shapes:triangles}{segitiga siku-siku}: hitunglah tinggi
itu, lalu luas segitiganya.
\end{exercise}

\begin{exercise}[$\star\star$]\label{exo:g8:pythagoras:7}
Layar televisi adalah persegi panjang $16{:}9$ berukuran $88.5$ cm
kali $49.8$ cm. Layar dijual menurut diagonalnya, dalam inci
($1$ inci $= 2.54$ cm). Hitunglah diagonalnya dalam cm, lalu dalam
inci (\omterm{def:g4:numbers:round}{bulatkan} ke inci terdekat).
\end{exercise}

\begin{exercise}[$\star\star$]\label{exo:g8:pythagoras:8}
Sebuah perahu berlayar $24$ km ke timur, lalu $10$ km ke utara.
Berapa jarak perahunya dari titik awalnya (lurus menyeberang)?
Gambarlah keadaannya lebih dahulu.
\end{exercise}

\begin{exercise}[$\star\star$]\label{exo:g8:pythagoras:9}
Tangga sepanjang $2.5$ m harus mencapai ambang jendela $2.4$ m di atas
tanah. Berapa jauh dari tembok kakinya harus diletakkan? Apakah
jawabannya cocok dengan anjuran keselamatan (kakinya kira-kira
\omterm{def:g3:division:half}{seperempat} panjang tangganya dari tembok)?
\end{exercise}

\begin{exercise}[$\star\star$]\label{exo:g8:pythagoras:10}
Pada kisi berkoordinat, gambarlah $A(1, 2)$ dan $B(5, 5)$, lalu
gambarlah \omterm{def:g6:shapes:triangles}{segitiga siku-siku} yang sisi siku-sikunya \omterm{def:g6:lines:perp}{sejajar} dengan
sumbunya dan sisi miringnya $[AB]$. Hitunglah $AB$. (Pelukisan ini,
untuk dua titik mana pun, menjadi rumus jarak pada geometri
berkoordinat, di buku Sekolah Menengah.)
\end{exercise}

\begin{exercise}[$\star\star\star$]\label{exo:g8:pythagoras:11}
Empat salinan \omterm{def:g6:shapes:triangles}{segitiga siku-siku} dengan sisi siku-siku $a$, $b$ dan
sisi miring $c$ diletakkan di pojok persegi bersisi $a + b$, sehingga
menyisakan persegi miring bersisi $c$ tanpa tertutupi di tengahnya.
\begin{enumerate}
  \item Nyatakan luas persegi besarnya dengan dua cara: secara
        langsung, dan sebagai (empat segitiga) $+$ (persegi miring).
  \item Jabarkan $(a + b)^2$ (lihat \cref{thm:g8:equations:expand})
        lalu simpulkan $c^2 = a^2 + b^2$: kamu sudah membuktikan
        teorema Pythagoras.
\end{enumerate}
\end{exercise}

\section{Soal: Tripel Pythagoras}

\begin{problem}[{Soal akhir pekan --- segitiga siku-siku berbilangan
bulat, dan kesamaan $(m^2 - n^2)^2 + (2mn)^2 = (m^2 + n^2)^2$ yang
membangkitkan semuanya}]
\label{pb:g8:pythagoras:1}
\emph{Tripel Pythagoras}\index{tripel Pythagoras} adalah tripel
bilangan bulat $(a, b, c)$ dengan
\[
a^2 + b^2 = c^2 ,
\]
seperti $(3, 4, 5)$ milik tukang bangunan. Sebuah lempeng tanah liat
Babilonia mendaftar tripel yang raksasa --- $(4601,\ 4800,\ 6649)$ di
antaranya --- lebih dari seribu tahun sebelum Pythagoras. Soal ini
memburu tripel, lalu membangun mesin yang menghasilkannya: sebuah
kesamaan aljabar yang langsung datang dari \cref{ch:g8:equations},
yang dipekerjakan untuk melayani geometri.

\medskip
\noindent\textbf{Bagian I --- Memburu tripel.}
\begin{enumerate}
  \item Periksalah bahwa $(5, 12, 13)$, $(8, 15, 17)$ dan
        $(20, 21, 29)$ adalah \omterm{pb:g8:pythagoras:1}{tripel Pythagoras}. Apa yang dikatakan
        \cref{thm:g8:pythagoras:converse} tentang segitiga yang
        panjang sisinya demikian?
  \item Jelaskan mengapa tali tegang dengan dua belas ruas yang sama,
        yang ditata sebagai segitiga bersisi $3$, $4$ dan $5$ ruas,
        memberi tukang bangunan \omterm{def:g3:shapes:rightangle}{sudut siku-siku} yang sempurna.
  \item Tunjukkan bahwa kalau $(a, b, c)$ \omterm{pb:g8:pythagoras:1}{tripel Pythagoras}, maka
        $(ka, kb, kc)$ juga, untuk setiap bilangan bulat $k \geq 1$.
        Simpulkan tiga tripel baru dari $(3, 4, 5)$.
  \item \omterm{def:g5:division:multiple}{Kelipatan} $(3, 4, 5)$ adalah tripel
        $(3k, 4k, 5k)$. Tunjukkan bahwa $(5, 12, 13)$ \emph{bukan}
        salah satunya.
  \item Jelaskan mengapa pertanyaan 3 dan 4 bersama-sama membuktikan
        bahwa menskalakan $(3, 4, 5)$ tidak pernah dapat menghasilkan
        setiap \omterm{pb:g8:pythagoras:1}{tripel Pythagoras}: mesin yang lain diperlukan.
\end{enumerate}

\medskip
\noindent\textbf{Bagian II --- Mesinnya.}
Ambillah dua bilangan bulat $m > n \geq 1$ lalu bentuklah ketiga bilangan
\[
a = m^2 - n^2,
\qquad
b = 2mn,
\qquad
c = m^2 + n^2 .
\]
\begin{enumerate}[resume]
  \item Dengan memakai kesamaan istimewa
        \cref{pb:g8:equations:1} (tulislah $A = m^2$ dan $B = n^2$),
        jabarkan $a^2$ dan $b^2$, lalu buktikan bahwa
        \[
        \left(m^2 - n^2\right)^2 + (2mn)^2
        = \left(m^2 + n^2\right)^2 :
        \]
        mesinnya selalu mengeluarkan \omterm{pb:g8:pythagoras:1}{tripel Pythagoras}.
  \item Jalankan mesinnya pada $(m, n) = (2, 1)$, $(3, 1)$, $(3, 2)$,
        $(4, 1)$ dan $(4, 3)$, lalu tampilkan kelima tripelnya dalam
        sebuah tabel.
  \item Satu tripel pada tabelmu adalah salinan terskala dari tripel
        yang lebih kecil. Yang mana, dan dari yang mana?
  \item Carilah $(m, n)$ yang membuat mesinnya mengeluarkan tripel
        $(20, 21, 29)$ pada pertanyaan 1 (sisi siku-sikunya dalam
        urutan mana pun).
  \item Carilah $m$ dan $n$ dengan $m^2 + n^2 = 100$, lalu simpulkan
        \omterm{pb:g8:pythagoras:1}{tripel Pythagoras} yang sisi miringnya $100$.
\end{enumerate}

\medskip
\noindent\textbf{Bagian III --- Keluarga, dan penutup tentang keganjilan.}
\begin{enumerate}[resume]
  \item Jalankan mesinnya dengan $m = n + 1$. Tunjukkan bahwa sisi
        siku-siku yang \omterm{def:g3:numbers:evenodd}{ganjil} adalah $2n + 1$, yang \omterm{def:g3:numbers:evenodd}{genap}
        $2n^2 + 2n$, dan sisi miringnya
        $2n^2 + 2n + 1$ --- jadi sisi miring dan sisi siku-siku yang
        \omterm{def:g3:numbers:evenodd}{genap} berbeda tepat $1$. Tulislah tripelnya untuk
        $n = 1, 2, 3, 4$.
  \item Simpulkan bahwa \emph{setiap} \omterm{def:g3:numbers:evenodd}{bilangan ganjil}
        $3, 5, 7, 9, \dots$ adalah sisi suatu \omterm{def:g6:shapes:triangles}{segitiga siku-siku}
        berbilangan bulat, lalu sebutkan tripel yang memuat $11$.
  \item \omterm{def:g3:numbers:evenodd}{Bilangan genap} juga berhasil: jalankan mesinnya dengan
        $n = 1$, lalu hasilkan tripel yang memuat $14$. Lalu
        perolehlah tripel yang sama untuk kedua kalinya, dengan
        menskalakan salah satu tripel pada tabel pertanyaan 7.
  \item Tunjukkan bahwa kuadrat \omterm{def:g3:numbers:evenodd}{bilangan genap} itu \omterm{def:g3:numbers:evenodd}{genap} dan bahwa
        kuadrat \omterm{def:g3:numbers:evenodd}{bilangan ganjil} itu \omterm{def:g3:numbers:evenodd}{ganjil}. (Tulislah bilangannya
        sebagai $2k$ atau $2k + 1$ lalu pakailah sebuah kesamaan.)
  \item Simpulkan penutup tentang keganjilannya: tidak ada \omterm{pb:g8:pythagoras:1}{tripel
        Pythagoras} yang ketiga bilangannya \omterm{def:g3:numbers:evenodd}{ganjil} --- pada setiap
        \omterm{def:g6:shapes:triangles}{segitiga siku-siku} berbilangan bulat, sedikitnya satu sisinya
        \omterm{def:g3:numbers:evenodd}{genap}.

\end{enumerate}
\end{problem}
